\subsubsection{Ford-Fulkerson (DFS max flow)}
\begin{lstlisting}[style=mystyle, language=C++]
// TC: O(M * max_flow) con DFS
// usa: max flow clasico; busca caminos aumentantes con DFS
// diferencia con Edmonds-Karp: FF usa DFS (depende del flujo), EK usa BFS O(N * M^2)
#include <bits/stdc++.h>
using namespace std;

const long long INF = 1e18;

long long dfs(int node, int t, long long pushed, vector<vector<pair<int,long long>>>& adj, vector<int>& vis){
	if(node == t || pushed == 0) return pushed;
	vis[node] = 1;
	for(auto& [adjN, cap] : adj[node]){
		if(!vis[adjN] && cap > 0){
			long long tr = dfs(adjN, t, min(pushed, cap), adj, vis);
			if(tr > 0){
				cap -= tr;
				for(auto& [backNode, backCap] : adj[adjN]){
					if(backNode == node){
						backCap += tr;
						break;
					}
				}
				return tr;
			}
		}
	}
	return 0;
}

long long fordFulkerson(int n, vector<vector<pair<int,long long>>>& adj, int s, int t){
	long long flow = 0;
	while(true){
		vector<int> vis(n, 0);
		long long pushed = dfs(s, t, INF, adj, vis);
		if(pushed == 0) break; // no existe camino aumentante
		flow += pushed;
	}
	return flow;
}

int main(){
	// Grafo de 4 nodos (mismo ejemplo que Edmonds-Karp)
	vector<vector<pair<int,long long>>> adj(4);
	auto add = [&](int u, int v, long long c){
		adj[u].push_back({v, c});
		adj[v].push_back({u, 0});
	};
	add(0,1,10); add(0,2,5); add(1,2,15); add(1,3,10); add(2,3,10);
	cout << "max flow = " << fordFulkerson(4, adj, 0, 3) << "\n"; // 15
	return 0;
}
\end{lstlisting}
